% GENERATED FILE. Edit canonical lessons, then run npm run generate:books.
% canonical-source-hash: fnv1a64:72d0d89b404c739a
% canonical-lessons: KA-C533-kasuti, KA-C533-karakusala, KA-C533-hene, KA-C533-ponisu, KA-C533-singarisu

\chapter{Handwork}
\label{ch:ka-handwork}
\hlchaptermodality{533}

\begin{chapteropening}
\textbf{By the end of this chapter:} \emph{I can say embroidery, handicraft, to knit, to string beads and to adorn.}

Complete the last lesson of chapter 533: I can say embroidery, handicraft, to knit, to string beads and to adorn.
\end{chapteropening}

\section[\texorpdfstring{\kn{ಕಸೂತಿ} (kasūti)}{ಕಸೂತಿ (kasūti)}]{\kn{ಕಸೂತಿ} (kasūti) — embroidery}

\label{lesson:KA-C533-kasuti}



\begin{quote}
Before the new one: say the Kannada for a lullaby, then the Kannada for to reminisce.
\end{quote}

\subsection*{You'll want to know: \kn{ಕಸೂತಿ}}
\textbf{\kn{ಕಸೂತಿ}} — \emph{kasūti} — \textquotedblleft{}embroidery\textquotedblright{}.

\begin{quote}
\kn{ಅಮ್ಮ} \kn{ಕರವಸ್ತ್ರದ} \kn{ಮೇಲೆ} \kn{ಕಸೂತಿ} \kn{ಮಾಡಿದರು}.
\end{quote}

\begin{quote}
\emph{amma karavastrada mēle kasūti māḍidaru.}
\end{quote}

\begin{quote}
Mother embroidered a handkerchief.
\end{quote}

\kn{ಕಸೂತಿ} \kn{ಮಾಡು}, to embroider.

\begin{cousinweb}
\textbf{\kn{ಕಸೂತಿ}} (\emph{kasūti}) is the Kannada word for embroidery; its origin is not settled.
\end{cousinweb}

\begin{grammarlens}[title={What you've built}]
One more word for describing what you lived and felt.
\end{grammarlens}

\subsection*{Your turn}
\begin{itemize}
\raggedright
  \item \emph{Say it:} \emph{kasūti}
  \item \emph{Say it:} \emph{kasūti}, once more
  \item \emph{Say it:} \emph{amma karavastrada mēle kasūti māḍidaru.} — Mother embroidered a handkerchief.
  \item \emph{Recall:} say the Kannada for a lullaby, then the Kannada for to reminisce, then say \emph{kasūti} again
\end{itemize}

\begin{tcolorbox}[breakable,colback=teal!4,colframe=teal!35!black,title={Before you move on}]
What does \kn{ಕಸೂತಿ} mean? (\textquotedblleft{}Embroidery\textquotedblright{}.) Say it once more.
\end{tcolorbox}
\section[\texorpdfstring{\kn{ಕರಕುಶಲ} (karakuśala)}{ಕರಕುಶಲ (karakuśala)}]{\kn{ಕರಕುಶಲ} (karakuśala) — handicraft}

\label{lesson:KA-C533-karakusala}



\begin{quote}
Before the new one: say the Kannada for to reminisce, then the Kannada for embroidery.
\end{quote}

\subsection*{You'll want to know: \kn{ಕರಕುಶಲ}}
\textbf{\kn{ಕರಕುಶಲ}} — \emph{karakuśala} — \textquotedblleft{}handicraft\textquotedblright{}.

\begin{quote}
\kn{ಜಾತ್ರೆಯಲ್ಲಿ} \kn{ಕರಕುಶಲ} \kn{ವಸ್ತುಗಳನ್ನು} \kn{ಮಾರುತ್ತಿದ್ದರು}.
\end{quote}

\begin{quote}
\emph{jātreyalli karakuśala vastugaḷannu māruttiddaru.}
\end{quote}

\begin{quote}
Handicrafts were being sold at the fair.
\end{quote}

\kn{ಕರ}, hand + \kn{ಕುಶಲ}, skilful.

\begin{cousinweb}
\textbf{\kn{ಕರಕುಶಲ}} (\emph{karakuśala}) is a learned compound built from Sanskrit \emph{kara}, \textquotedblleft{}the hand\textquotedblright{}, and \emph{kuśala}, \textquotedblleft{}skilful, clever\textquotedblright{} (Monier-Williams).
\end{cousinweb}

\begin{grammarlens}[title={What you've built}]
One more word for describing what you lived and felt.
\end{grammarlens}

\subsection*{Your turn}
\begin{itemize}
\raggedright
  \item \emph{Say it:} \emph{karakuśala}
  \item \emph{Say it:} \emph{karakuśala}, once more
  \item \emph{Say it:} \emph{jātreyalli karakuśala vastugaḷannu māruttiddaru.} — Handicrafts were being sold at the fair.
  \item \emph{Recall:} say the Kannada for to reminisce, then the Kannada for embroidery, then say \emph{karakuśala} again
\end{itemize}

\begin{tcolorbox}[breakable,colback=teal!4,colframe=teal!35!black,title={Before you move on}]
What does \kn{ಕರಕುಶಲ} mean? (\textquotedblleft{}Handicraft\textquotedblright{}.) Say it once more.
\end{tcolorbox}
\section[\texorpdfstring{\kn{ಹೆಣೆ} (heṇe)}{ಹೆಣೆ (heṇe)}]{\kn{ಹೆಣೆ} (heṇe) — to knit; to plait}

\label{lesson:KA-C533-hene}



\begin{quote}
Before the new one: say the Kannada for embroidery, then the Kannada for handicraft.
\end{quote}

\subsection*{You'll want to know: \kn{ಹೆಣೆ}}
\textbf{\kn{ಹೆಣೆ}} — \emph{heṇe} — \textquotedblleft{}to knit; to plait\textquotedblright{}.

\begin{quote}
\kn{ಅಜ್ಜಿ} \kn{ಚಳಿಗಾಲಕ್ಕೆ} \kn{ಉಣ್ಣೆಯ} \kn{ಸ್ವೆಟರ್} \kn{ಹೆಣೆದರು}.
\end{quote}

\begin{quote}
\emph{ajji caḷigālakke uṇṇeya sveṭar heṇedaru.}
\end{quote}

\begin{quote}
Grandmother knitted a woollen sweater for the winter.
\end{quote}

Past (I …): \kn{ಹೆಣೆದೆ}, \emph{heṇede}.

Hair too: \kn{ಜಡೆ} \kn{ಹೆಣೆ}, to plait a braid.

\begin{cousinweb}
\textbf{\kn{ಹೆಣೆ}} (\emph{heṇe}) is native Dravidian: DEDR 4160 lists Kannada \emph{peṇe}, \textquotedblleft{}to tie together, intertwine, twist, plait, braid\textquotedblright{}. Old p- became h-.
\end{cousinweb}

\begin{grammarlens}[title={What you've built}]
One more word for describing what you lived and felt.
\end{grammarlens}

\subsection*{Your turn}
\begin{itemize}
\raggedright
  \item \emph{Say it:} \emph{heṇe}
  \item \emph{Say it:} \emph{heṇe}, once more
  \item \emph{Say it:} \emph{ajji caḷigālakke uṇṇeya sveṭar heṇedaru.} — Grandmother knitted a woollen sweater for the winter.
  \item \emph{Recall:} say the Kannada for embroidery, then the Kannada for handicraft, then say \emph{heṇe} again
\end{itemize}

\begin{tcolorbox}[breakable,colback=teal!4,colframe=teal!35!black,title={Before you move on}]
What does \kn{ಹೆಣೆ} mean? (\textquotedblleft{}To knit; to plait\textquotedblright{}.) Say it once more.
\end{tcolorbox}
\section[\texorpdfstring{\kn{ಪೋಣಿಸು} (pōṇisu)}{ಪೋಣಿಸು (pōṇisu)}]{\kn{ಪೋಣಿಸು} (pōṇisu) — to string beads}

\label{lesson:KA-C533-ponisu}



\begin{quote}
Before the new one: say the Kannada for handicraft, then the Kannada for to knit.
\end{quote}

\subsection*{You'll want to know: \kn{ಪೋಣಿಸು}}
\textbf{\kn{ಪೋಣಿಸು}} — \emph{pōṇisu} — \textquotedblleft{}to string beads\textquotedblright{}.

\begin{quote}
\kn{ತಂಗಿ} \kn{ಮಲ್ಲಿಗೆ} \kn{ಹೂವುಗಳನ್ನು} \kn{ದಾರದಲ್ಲಿ} \kn{ಪೋಣಿಸಿದಳು}.
\end{quote}

\begin{quote}
\emph{taṅgi mallige hūvugaḷannu dāradalli pōṇisidaḷu.}
\end{quote}

\begin{quote}
My younger sister strung jasmine flowers on a thread.
\end{quote}

Past (I …): \kn{ಪೋಣಿಸಿದೆ}, \emph{pōṇiside}.

Beads, pearls or flowers on a thread.

\begin{cousinweb}
\textbf{\kn{ಪೋಣಿಸು}} (\emph{pōṇisu}) is native Dravidian: DEDR 4584 lists Kannada \emph{pōṇisu}, \emph{pavaṇisu}, \textquotedblleft{}to couple, to string together, to thread\textquotedblright{}, beside Telugu \emph{pōhaṇincu}.
\end{cousinweb}

\begin{grammarlens}[title={What you've built}]
One more word for describing what you lived and felt.
\end{grammarlens}

\subsection*{Your turn}
\begin{itemize}
\raggedright
  \item \emph{Say it:} \emph{pōṇisu}
  \item \emph{Say it:} \emph{pōṇisu}, once more
  \item \emph{Say it:} \emph{taṅgi mallige hūvugaḷannu dāradalli pōṇisidaḷu.} — My younger sister strung jasmine flowers on a thread.
  \item \emph{Recall:} say the Kannada for handicraft, then the Kannada for to knit, then say \emph{pōṇisu} again
\end{itemize}

\begin{tcolorbox}[breakable,colback=teal!4,colframe=teal!35!black,title={Before you move on}]
What does \kn{ಪೋಣಿಸು} mean? (\textquotedblleft{}To string beads\textquotedblright{}.) Say it once more.
\end{tcolorbox}
\section[\texorpdfstring{\kn{ಸಿಂಗರಿಸು} (siṅgarisu)}{ಸಿಂಗರಿಸು (siṅgarisu)}]{\kn{ಸಿಂಗರಿಸು} (siṅgarisu) — to adorn; to dress up}

\label{lesson:KA-C533-singarisu}



\begin{quote}
Before the new one: say the Kannada for to knit, then the Kannada for to string beads.
\end{quote}

\subsection*{You'll want to know: \kn{ಸಿಂಗರಿಸು}}
\textbf{\kn{ಸಿಂಗರಿಸು}} — \emph{siṅgarisu} — \textquotedblleft{}to adorn; to dress up\textquotedblright{}.

\begin{quote}
\kn{ಹಬ್ಬಕ್ಕೆ} \kn{ಮನೆಯನ್ನು} \kn{ಹೂವುಗಳಿಂದ} \kn{ಸಿಂಗರಿಸಿದರು}.
\end{quote}

\begin{quote}
\emph{habbakke maneyannu hūvugaḷinda siṅgarisidaru.}
\end{quote}

\begin{quote}
They decorated the house with flowers for the festival.
\end{quote}

Past (I …): \kn{ಸಿಂಗರಿಸಿದೆ}, \emph{siṅgariside}.

Close to \kn{ಅಲಂಕರಿಸು}; \kn{ಸಿಂಗಾರ} is finery.

\begin{cousinweb}
\textbf{\kn{ಸಿಂಗರಿಸು}} (\emph{siṅgarisu}) is built on \emph{siṅgāra}, \textquotedblleft{}finery\textquotedblright{}, a Kannada form of Sanskrit \emph{śṛṅgāra}, which Monier-Williams glosses as love and also \textquotedblleft{}elegant dress, finery\textquotedblright{} (Monier-Williams).
\end{cousinweb}

\begin{grammarlens}[title={What you've built}]
One more word for describing what you lived and felt.
\end{grammarlens}

\subsection*{Your turn}
\begin{itemize}
\raggedright
  \item \emph{Say it:} \emph{siṅgarisu}
  \item \emph{Say it:} \emph{siṅgarisu}, once more
  \item \emph{Say it:} \emph{habbakke maneyannu hūvugaḷinda siṅgarisidaru.} — They decorated the house with flowers for the festival.
  \item \emph{Recall:} say the Kannada for to knit, then the Kannada for to string beads, then say \emph{siṅgarisu} again
\end{itemize}

\begin{tcolorbox}[breakable,colback=teal!4,colframe=teal!35!black,title={Before you move on}]
What does \kn{ಸಿಂಗರಿಸು} mean? (\textquotedblleft{}To adorn; to dress up\textquotedblright{}.) Say it once more.
\end{tcolorbox}
